\makeatletter
\usepackage{mathtools,bm,microtype,calc,array,pdfpages,nicematrix,blkarray,tikz,tikz-cd,booktabs,hyphenat,csquotes,hyperref,cleveref, siunitx, xcolor, pifont, bbding}
\RequirePackage{biblatex}
\usepackage[flushmargin]{footmisc} % Remove footnote indentation
\graphicspath{{visualizations/}}
\renewbibmacro{in:}{, }
\DeclareFieldFormat[misc]{title}{\mkbibquote{#1}}
%\DeclareFieldFormat[article]{volume}{\mkbibbold{#1}}
%\DeclareFieldFormat{url}{\textsc{url}: \href{#1}{#1}}
%\DeclareFieldFormat{doi}{\textsc{doi}: \href{http://dx.doi.org/#1}{#1}}
%\DeclareFieldFormat{eprint:arxiv}{arXiv:\href{https://arxiv.org/abs/#1}{#1}}
\addbibresource{biblio.bib}

%\usepackage{hyperref,cleveref}

% This patch accounts for a bug recently introduced in latex, see https://tex.stackexchange.com/q/730148/27717
\newcommand\patchtheorem[2]{\AddToHook{env/#1/begin}{\crefalias{#2}{#1}}}
\patchtheorem{lemma}{theorem}
\patchtheorem{proposition}{theorem}
\patchtheorem{corollary}{theorem}
\patchtheorem{conjecture}{theorem}
\patchtheorem{assumption}{theorem}
\patchtheorem{observation}{theorem}
\patchtheorem{example}{theorem}
\patchtheorem{definition}{theorem}
\patchtheorem{remark}{theorem}
\patchtheorem{question}{theorem}
\patchtheorem{problem}{theorem}

% Enable sub-theorems and references to them
\usepackage[patch-newtheorem]{mytheorem-enum}

\crefname{equation}{}{}
\Crefname{page}{page}{Pages}

\theoremstyle{plain}

\theoremstyle{remark}
\newtheorem{assumption}[theorem]{Assumption}
\newtheorem{observation}[theorem]{Observation}

\newtheorem{claim}{Claim}[theorem]
\renewcommand{\theclaim}{\arabic{claim}}
\newtheorem*{claim*}{Claim}

\newenvironment{claimproof}[1][\proofname\ of Claim]{%
	\par\normalfont\topsep6\p@\@plus6\p@\relax
	\trivlist
	\item[\hskip\labelsep\itshape#1\@addpunct{.}]\ignorespaces
}{\endtrivlist\@endpefalse}


\Crefname{example}{Example}{Examples}
\Crefname{theorem}{Theorem}{Theorems}
\Crefname{lemma}{Lemma}{Lemmas}
\Crefname{proposition}{Proposition}{Propositions}
\Crefname{figure}{Figure}{Figures}
\Crefname{remark}{Remark}{Remarks}
\Crefname{corollary}{Corollary}{Corollaries}
\Crefname{conjecture}{Conjecture}{Conjectures}

\usepackage{xspace}
\newcommand\OSCAR{\texttt{OSCAR}\xspace}
\newcommand\polymake{\texttt{polymake}\xspace}
\newcommand\Macaulay{\texttt{Macaulay2}\xspace}

\newcommand\CC{{\mathbb C}}
\newcommand\EE{{\mathbb E}}
\newcommand\LL{{\mathbb L}}
\newcommand\NN{{\mathbb N}}
\newcommand\KK{{\mathbb K}}
\newcommand\PP{{\mathbb P}}
\newcommand\QQ{{\mathbb Q}}
\newcommand\RR{{\mathbb R}}
\newcommand\Sph{{\mathbb S}}
\newcommand\TT{{\mathbb T}}
\newcommand\ZZ{{\mathbb Z}}
\newcommand\FF{{\mathbb F}}
\newcommand\RP{\mathbb{RP}}

\newcommand\cD{{\mathcal D}}
\DeclareMathOperator{\Gin}{gin}
\DeclareMathOperator{\In}{in}
\newcommand\frM{{\mathfrak m}}
\newcommand\frR{{\mathfrak r}}
\newcommand\frU{{\mathfrak u}}
\newcommand\frV{{\mathfrak v}} % was previously defined in the middle of the main tex
\newcommand\frX{{\mathfrak x}}
\newcommand\frY{\mathfrak{y}}
\let\xto\xrightarrow
\let\symdiff\bigtriangleup
\newcommand{\Matr}{\mathcal{M}}
\DeclareMathOperator{\argmin}{arg\,min}
\DeclareMathOperator{\rk}{rk}
\newcommand\GF[1]{\operatorname{GF}(#1)}
\newcommand{\GL}{\mathrm{GL}}
\newcommand{\PSL}{\mathrm{PSL}}
\DeclareMathOperator{\Gr}{Gr}
\DeclarePairedDelimiterXPP{\SymmetricGroup}[1]{\operatorname{Sym}}{(}{)}{}{#1}
\newcommand{\PG}{\mathrm{PG}}
\DeclareMathOperator{\inv}{Inv}
\DeclareMathOperator{\card}{card}
\newcommand{\RepVert}[3]{#1|_{#2}^{#3}}
\newcommand{\NinLSpan}[2]{\mathrm{LLI}(#1, #2)}
\newcommand{\lex}[1]{#1_{\text{lex}}}
\newcommand{\InitialSegment}{\mathrm{In}}
\newcommand{\ShiftGraph}{\mathrm{PSG}}
\newcommand{\ContractedGraph}{\overline{\ShiftGraph}}
\newcommand{\CShift}{\Gamma}
\DeclareMathOperator{\Char}{char}

\newcommand{\actson}{\mathbin{\rotatebox[origin=c]{-90}{$\circlearrowleft$}}}
\newcommand{\actsfrom}{\mathbin{\rotatebox[origin=c]{90}{$\circlearrowright$}}}

% Set builder notation. Call by \Set{v; P(v)}, \Set*{v; P(v)} for auto-sized delimiters, or \Set[\big]{v; P(v)} for fixed size.
\DeclarePairedDelimiterX{\Set}[1]{\{}{\}}{\setargs{#1}}

% Linear span. Like \Set. Optionally takes the field as subscript; e.g., \Span_{k}*{v} or \Span_{k}[\big]{v}$.
\NewDocumentCommand{\Span}{E{_}{{}}}{\operatorname{span}_{#1}\SpanX}

% Use as \abs{}, \abs*{} or \abs[\big]{}
\DeclarePairedDelimiter{\abs}{\lvert}{\rvert}
\DeclarePairedDelimiter{\floor}{\lfloor}{\rfloor}
\DeclarePairedDelimiter{\ceil}{\lceil}{\rceil}
% Internals for the above
\DeclarePairedDelimiterX{\SpanX}[1]{(}{)}{\setargs{#1}}
\NewDocumentCommand{\setargs}{>{\SplitArgument{1}{;}}m}{\setargsaux#1}
%\NewDocumentCommand{\setargsaux}{mm}{\IfNoValueTF{#2}{#1}{#1\nonscript\:\delimsize\vert\allowbreak\nonscript\:\mathopen{}#2}}%
\NewDocumentCommand{\setargsaux}{mm}{\IfNoValueTF{#2}{#1}{#1:\allowbreak #2}}%

% Small-script case distinction environment.
\newenvironment{smallcases}{\left\{\begin{smallmatrix*}[l]}{\end{smallmatrix*}\right.}

% Cases with brace on both sides (https://tex.stackexchange.com/a/308762/27717)
\newcases{lrcases}{\quad}
{$\m@th\textstyle{##}$\hfil}
{$\m@th\textstyle{##}$\hfil}
{\lbrace}
{\rbrace}
\newcases{lrdcases*}
{\quad}
{$\m@th\textstyle{##}$\hfil}
{{##}\hfil}
{\lbrace}
{\rbrace}

% Align in substack (https://tex.stackexchange.com/a/198806/27717)
\newcommand{\subalign}[1]{%
	\vcenter{%
		\Let@ \restore@math@cr \default@tag
		\baselineskip\fontdimen10 \scriptfont\tw@
		\advance\baselineskip\fontdimen12 \scriptfont\tw@
		\lineskip\thr@@\fontdimen8 \scriptfont\thr@@
		\lineskiplimit\lineskip
		\ialign{\hfil$\m@th\scriptstyle##$&$\m@th\scriptstyle{}##$\hfil\crcr
			#1\crcr
		}%
	}%
}

% Produce command \MakeCommaBreakable that makes comma breakable in inline math (https://tex.stackexchange.com/a/716369/27717.
{
	\gdef\OldComma{,}
	\catcode`\,=13
	\gdef\MakeCommaBreakable{%
		\catcode`\,=13
		\def,{%
			\ifmmode
			\OldComma\>\allowbreak
			\else
			\OldComma
			\fi
		}%
	}
}

% Small version of array (https://tex.stackexchange.com/a/98692/27717)
\newenvironment{smallarray}[1]{%
	\null\,\vcenter\bgroup\scriptsize%
	\renewcommand{\arraystretch}{0.7}%
	\arraycolsep=.13885em%
	\hbox\bgroup$\array{@{}#1@{}}%
}{%
	\endarray$\egroup\egroup\,\null%
}

% Scaling version of \vdots and \ddots (better than mathdots package, see https://tex.stackexchange.com/a/246695/27717)
\usepackage{letltxmacro}
\LetLtxMacro\orgvdots\vdots
\LetLtxMacro\orgddots\ddots
\DeclareRobustCommand\vdots{%
	\mathpalette\@vdots{}%
}
\newcommand*{\@vdots}[2]{%
	% #1: math style
	% #2: unused
	\sbox0{$#1\cdotp\cdotp\cdotp\m@th$}%
	\sbox2{$#1.\m@th$}%
	\vbox{%
		\dimen@=\wd0 %
		\advance\dimen@ -3\ht2 %
		\kern.5\dimen@
		% remove side bearings
		\dimen@=\wd2 %
		\advance\dimen@ -\ht2 %
		\dimen2=\wd0 %
		\advance\dimen2 -\dimen@
		\vbox to \dimen2{%
			\offinterlineskip
			\copy2 \vfill\copy2 \vfill\copy2 %
		}%
	}%
}
\DeclareRobustCommand\ddots{%
	\mathinner{%
		\mathpalette\@ddots{}%
		\mkern\thinmuskip
	}%
}
\newcommand*{\@ddots}[2]{%
	% #1: math style
	% #2: unused
	\sbox0{$#1\cdotp\cdotp\cdotp\m@th$}%
	\sbox2{$#1.\m@th$}%
	\vbox{%
		\dimen@=\wd0 %
		\advance\dimen@ -3\ht2 %
		\kern.5\dimen@
		% remove side bearings
		\dimen@=\wd2 %
		\advance\dimen@ -\ht2 %
		\dimen2=\wd0 %
		\advance\dimen2 -\dimen@
		\vbox to \dimen2{%
			\offinterlineskip
			\hbox{$#1\mathpunct{.}\m@th$}%
			\vfill
			\hbox{$#1\mathpunct{\kern\wd2}\mathpunct{.}\m@th$}%
			\vfill
			\hbox{$#1\mathpunct{\kern\wd2}\mathpunct{\kern\wd2}\mathpunct{.}\m@th$}%
		}%
	}%
}

% %%%%%%%%%%%%%%%%%%%%%%%%%%%%%%%%%%%%%%%%%%%%%%%%%%%%%%%%%%%%%%%%%%%%%
%% % % preliminary version only
%\usepackage[colorinlistoftodos,bordercolor=orange,backgroundcolor=orange!20,linecolor=orange,textsize=tiny]{todonotes}
%\usepackage{prelim2e}
%\renewcommand{\PrelimText}{\raisebox{-\height}[0pt][0pt]{\sf\scriptsize --- INCOMPLETE DRAFT of \today ---}}
%\setlength{\marginparwidth}{1.2in}

%% % %%%%%%%%%%%%%%%%%%%%%%%%%%%%%%%%%%%%%%%%%%%%%%%%%%%%%%%%%%%%%%%%%%%%%

% Make \text{\todo{...}} work in math; see https://tex.stackexchange.com/a/132451/27717
% Michael: I do not like it, because this code does not correctly separate overlapping boxes
%\usepackage{marginnote}
%\makeatletter
% \renewcommand{\@todonotes@drawMarginNoteWithLine}{%
	% 	\begin{tikzpicture}[remember picture, overlay, baseline=-0.75ex]%
		% 		\node [coordinate] (inText) {};%
		% 	\end{tikzpicture}%
	% 	\marginnote[{% Draw note in left margin
		% 		\@todonotes@drawMarginNote%
		% 		\@todonotes@drawLineToLeftMargin%
		% 	}]{% Draw note in right margin
		% 		\@todonotes@drawMarginNote%
		% 		\@todonotes@drawLineToRightMargin%
		% 	}%
	% }

\newlength\ExtraCaptionMargin

% Footnote without number
\newcommand\blfootnote[1]{%
	\begingroup
	\renewcommand\thefootnote{}\footnote{#1}%
	\addtocounter{footnote}{-1}%
	\endgroup
}

\newcommand{\Email}[1]{\href{mailto:#1}{#1}}

\let\origDdots\ddots
\renewcommand{\ddots}[1][0]{\mathpalette\ddotsI{#1}}
\newcommand{\ddotsI}[2]{\rotatebox[origin=c]{#2}{$#1\origDdots$}}

\newcommand\SwapBelowDisplaySkip{\vskip-\belowdisplayskip\vskip\belowdisplayshortskip\noindent}

% Separates multiple citations by semicolon, rather than comma
\renewcommand\multicitedelim{\addsemicolon\space}

% From https://tex.stackexchange.com/a/19100
\mathchardef\breakingcomma\mathcode`\,
{\catcode`,=\active
	\gdef,{\breakingcomma\discretionary{}{}{}}
}
\newcommand{\mathlist}[1]{$\mathcode`\,=\string"8000 #1$}

\newcommand\Rot[1]{\rotatebox{90}{#1}}%
\newcommand\CS[1]{\csname #1\endcsname}%
\newcommand\Tick[1]{\resizebox{!}{\heightof{t}}{\ifthenelse{#1=1}{\Checkmark}{\XSolidBrush}}}%
\newcommand\EarlyReturn[1]{\ifthenelse{\equal{\CS{#1LvmaxLenBefore}}{n/a}}{*}{}}%
\newcommand\TimeOrDash[2][]{%
	\IfSubStr{\CS{#2Trials}}{>}%
	{{--}}%
	{\tablenum[table-auto-round,table-format=#1]{\CS{#2Time}}}%
}
\newcommand\FailedShift[1]{\IfSubStr{\CS{#1LvTrials}}{>}{${}^\dagger$}{}}%
\DeclareSIUnit{\noop}{\kern 0pt}%
\newcommand{\Mem}[1]{%
	\IfInteger{#1}{%
		\makebox[\widthof{\SI{10}{\giga\byte}}][r]{\SI[exponent-mode=engineering, exponent-to-prefix = true, round-mode=places, round-precision=0]{#1}{\noop}}
	}{#1}%
}%

% Create a new columntype A[fmt] := S[table-format=fmt]
\makeatletter
\newtoks\mytoks
\mytoks\expandafter{\the\NC@list}
\newcolumntype{A}{}
\NC@list\expandafter{\the\mytoks \NC@do A}
\renewcommand*{\NC@rewrite@A}[1][]{%
	\@temptokena\expandafter{%
		\the\@temptokena
		S[table-format=#1]
	}%
	\NC@find
}
\makeatother

% Create a new columntype M[fmt] := S[table-format=fmt]
\makeatletter
\newtoks\mytoks
\mytoks\expandafter{\the\NC@list}
\newcolumntype{M}{}
\NC@list\expandafter{\the\mytoks \NC@do M}
\renewcommand*{\NC@rewrite@M}[1][]{%
	\@temptokena\expandafter{%
		\the\@temptokena
		S[round-mode=places, round-precision=0, fixed-exponent=9, scientific-notation=fixed, drop-exponent, table-format=#1]
	}%
	\NC@find
}


\makeatother

